% ===== BIBLIOGRAPHIE (toute fin du document) =====
\clearpage\phantomsection\section*{Bibliographie}\addcontentsline{toc}{section}{Bibliographie}

Cette bibliographie liste toutes les sources citées dans le rapport : celles de la partie 2.1 (sources A à G) et celles de l'étude des besoins des plantes (parties 3.6 et 3.7, sources C, H, J, K, L, M et N). Je les ai retenues pour leur sérieux : articles de revues scientifiques, publications d'universités et d'organismes publics. Les chiffres du texte viennent de ces documents, et les formulations sont les miennes. Les liens sont cliquables, et les sources sont proposées pour validation. Je n'ai pas lu toutes les pages en entier : je m'appuie sur les passages qui donnent les chiffres cités.

\textbf{Source A.} Revue \emph{HortScience}, volume 59, numéro 8, pages 1202 et suivantes, American Society for Horticultural Science. Article d'histoire de l'hydroponie : Woodward en 1699, Knop en 1868, Gericke et le mot « hydroponie » en 1937. Lien : \url{https://journals.ashs.org/hortsci/}

\textbf{Source B.} USDA, National Agricultural Library, page « Hydroponics ». Rappelle Woodward (1699), Sachs et Knop, Gericke dans les années 1920 et la solution de Hoagland et Arnon en 1938. Lien : \url{https://nal.usda.gov/legacy/afsic/hydroponics}

\textbf{Source C.} Barbosa, G. L. et coauteurs (Arizona State University), « Comparison of Land, Water, and Energy Requirements of Lettuce Grown Using Hydroponic vs. Conventional Agricultural Methods », \emph{International Journal of Environmental Research and Public Health}, 2015, volume 12, pages 6879 et suivantes. DOI : 10.3390/ijerph120606879. Lien : \url{https://www.ncbi.nlm.nih.gov/pmc/articles/PMC4483736/}

\textbf{Source D.} Article de synthèse « Pathogenic and beneficial microorganisms in soilless cultures », \emph{Agronomy for Sustainable Development}, 2011, volume 31, pages 191 à 203. DOI : 10.1051/agro/2010018 (auteurs à confirmer avant citation définitive).

\textbf{Source E.} Revue sur les bactéries bénéfiques et pathogènes des systèmes hydroponiques, base PubMed Central, identifiant PMC11548230. Lien : \url{https://pmc.ncbi.nlm.nih.gov/articles/PMC11548230}

\textbf{Source F.} Wageningen University \& Research, document « Disinfection in hydroponic systems and hygiene ». Lien : \url{https://edepot.wur.nl/403803}

\textbf{Source G.} Gold, S. E. et Stanghellini, M. E., « Effects of temperature on Pythium root rot of spinach grown under hydroponic conditions », \emph{Phytopathology}, 1985, volume 75, pages 333 à 337. Cité dans la bibliographie de la source D.

\textbf{Source H.} Wheeler, R. M. et coauteurs, « Growth and Gas Exchange by Lettuce Stands in a Closed, Controlled Environment », \emph{Journal of the American Society for Horticultural Science}, 1994, 119(3), 610 à 615 : \textcolor{blue}{\url{https://journals.ashs.org/downloadpdf/view/journals/jashs/119/3/article-p610.pdf}}.

\textbf{Source J.} Kratky, B. A., « Three non-circulating hydroponic methods for growing lettuce », \emph{Acta Horticulturae}, 2009, 843, 65 à 72, université d'Hawaï : \textcolor{blue}{\url{https://www.ctahr.hawaii.edu/Hawaii/downloads/Three_Non-circulating_Hydroponic_Methods_for_Growing_lettuce.pdf}}.

\textbf{Source K.} Alabama Cooperative Extension System, « Greenhouse Lettuce Production », fiche ANR-2903 : \textcolor{blue}{\url{https://www.aces.edu/wp-content/uploads/2022/08/ANR-2903_GreenhouseLettuceProduction_072822L-G.pdf}}.

\textbf{Source L.} « Curva de absorção de nutrientes em alface hidropônica », \emph{Revista Caatinga}, janvier 2009 : \textcolor{blue}{\url{https://doaj.org/article/3cad1b7f657c4b9897b2eca48a74d30d}}.

\textbf{Source M.} Sharkey, A. et coauteurs, « Characterizing the Temporally Dynamic Nature of Relative Growth Rates: A Kinetic Analysis on Nitrogen-, Phosphorus-, and Potassium-Limited Growth », \emph{Agriculture}, juillet 2025 : \textcolor{blue}{\url{https://doaj.org/article/3f4baf90f7064b70be9ff31b0bc3bfbc}}.

\textbf{Source N.} Vought, K., Bayabil, H. et Martin-Ryals, A., « How Efficient Are Hydroponic Lettuce Production Systems in Nutrient Usage Under Low and High Salinity? », université de Floride (UF/IFAS), fiche AE610 : \textcolor{blue}{\url{https://edis.ifas.ufl.edu/publication/AE610/pdf}}.

\textbf{Limites.} L'étude de la source C repose sur un modèle de serre dans une région désertique des États-Unis, pas sur une mesure sur une tour comme la mienne. Les sources D à G n'ont pas toutes été lues en entier : je m'appuie sur leurs résumés. Je propose de valider ce point avec mon tuteur avant la version finale.
